% Part: proof-theory
% Chapter: sequent-calculus
% Section: invertibility
%READERNOTE{SOL6-A331: व्युत्क्रमण व आकुंचनाच्या आधारप्रसंगांत असत्याचे स्वयंसिद्धक समाविष्ट केले; प्रतिस्थापनापूर्वी अंतर्गत आयगेन स्थिरांक सुरक्षित ठेवले. उंची आणि आधारविधानांची संख्या वेगळ्या अक्षरांनी लिहिली.}

\documentclass[../../../include/open-logic-section]{subfiles}

\begin{document}

\olfileid{pt}{seq}{inv}

\olsection{नियमांची व्युत्क्रमणीयता}

\begin{defn}
एखादा नियम~$R$,
\[
\AxiomC{$S_1$}
\AxiomC{$\dots$}
\AxiomC{$S_n$}
\RightLabel{$R$}
\TrinaryInfC{$S$}
\DisplayProof
\]
पद्धतीत \emph{व्युत्क्रमणीय} असतो, जर $\Proves S$
असताना $i = 1$, \dots,~$n$ साठी $\Proves[h_i] S_i$
मिळत असेल. तो \emph{!!{height} जपून व्युत्क्रमणीय}
(किंवा \emph{!!{hp}-व्युत्क्रमणीय}) असतो, जर
$\Proves_h S$ असताना $i = 1$, \dots,~$n$
साठी $\Proves[h] S_i$ मिळत असेल.
\end{defn}

\begin{lem}\ollabel{lem:G3c-invert}
\Log{G3c} मधील $\lnot$, $\land$, $\lor$, आणि $\lif$
यांचे नियम !!{height} जपून व्युत्क्रमणीय आहेत;
म्हणजे:
\begin{enumerate}
  \item \RightR{\lnot}: $\Log{G3c} \Proves[n] \Gamma \Sequent
  \Delta, \lnot !A$ असेल, तर $\Log{G3c} \Proves[n] !A, \Gamma \Sequent
  \Delta$.
  \item \LeftR{\lnot}: $\Log{G3c} \Proves[n] \Gamma, \lnot !A
  \Sequent \Delta$ असेल, तर $\Log{G3c} \Proves[n] \Gamma \Sequent \Delta,
  !A$.
  \item \RightR{\land}: $\Log{G3c} \Proves[n] \Gamma \Sequent
  \Delta, !A \land !B$ असेल, तर $\Log{G3c} \Proves[n] \Gamma \Sequent
  \Delta, !A$ आणि $\Log{G3c} \Proves[n] \Gamma \Sequent
  \Delta, !B$.
  \item \LeftR{\land}: $\Log{G3c} \Proves[n] \Gamma, !A \land !B
  \Sequent \Delta$ असेल, तर $\Log{G3c} \Proves[n] !A, !B, \Gamma \Sequent
  \Delta$.
  \item \RightR{\lor}: $\Log{G3c} \Proves[n] \Gamma \Sequent
  \Delta, !A \lor !B$ असेल, तर $\Log{G3c} \Proves[n] \Gamma \Sequent
  \Delta, !A, !B$.
  \item \LeftR{\lor}: $\Log{G3c} \Proves[n] !A \lor !B, \Gamma \Sequent
  \Delta$ असेल, तर $\Log{G3c} \Proves[n] !A, \Gamma \Sequent
  \Delta$ आणि $\Log{G3c} \Proves[n] !B, \Gamma \Sequent
  \Delta$.
  \item \RightR{\lif}: $\Log{G3c} \Proves[n] \Gamma \Sequent
  \Delta, !A \lif !B$ असेल, तर $\Log{G3c} \Proves[n] !A, \Gamma \Sequent
  \Delta, !B$.
  \item \LeftR{\lif}: $\Log{G3c} \Proves[n] !A \lif !B, \Gamma
  \Sequent \Delta$ असेल, तर $\Log{G3c} \Proves[n] \Gamma \Sequent \Delta,
  !A$ आणि $\Log{G3c} \Proves[n] !B, \Gamma \Sequent \Delta$.
\end{enumerate}
\end{lem}

\begin{proof}
वरील प्रमेयिकेत दिलेल्या विधानीय नियम~$R$ साठी पुढील
गोष्ट दाखवायची आहे:~$R$ च्या निष्कर्ष~$S$ ची
!!a{proof}~$\pi$ असेल, तर $R$ च्या पूर्वपक्ष~$S_i$
यांच्या !!{proof}s~$\pi_i$ सुद्धा असतात, आणि
$\pheight{\pi_i}
\le \pheight{\pi}$. नियम~\LeftR{\land} विचारात घ्या:
अंतिम क्रमवर्ती $!A \land !B, \Gamma
\Sequent \Delta$ या रूपाची असलेली !!a{proof}~$\pi$
आहे असे समजा. $!A, !B, \Gamma \Sequent \Delta$
ची अशी !!a{proof}~$\pi'$ आहे की $\pheight{\pi'} \le
\pheight{\pi}$, हे दाखवायचे आहे. $n = \pheight{\pi}$
वर विगमन करून हे सिद्ध करू.

$n = 0$ असेल तर $\pi$ फक्त एकच स्वयंसिद्धक
असतो: $!A \land !B, \Gamma
\Sequent \Delta$. \Log{G3c} मधील स्वयंसिद्धके
$!C, \Gamma
\Sequent \Delta, !C$ या रूपाची असतात, आणि~$!C$
अण्वीय असतो. म्हणजेच $!C$ हा $!A \land !B$
असू शकत नाही. $!A \land !B, \Gamma \Sequent \Delta$
स्वयंसिद्धक असेल, तर एखादा अण्वीय~$!C$ हा~$\Gamma$
आणि~$\Delta$ दोहोंचा !!a{element} असला पाहिजे.
मग $!A, !B, \Gamma \Sequent \Delta$ हाही
स्वयंसिद्धक असतो आणि आपल्याला हवी असलेली
!!{proof}~$\pi'$ मिळाली (तिची !!{height} ही~$0$).
दुसरा स्वयंसिद्धक-प्रसंग म्हणजे~$\lfalse\in\Gamma$.
संयोगाची आस्थिती बदलल्यावरही तो डावीकडे राहतो;
म्हणून अपेक्षित क्रमवर्ती पुन्हा उंची~$0$ चे
असत्य-स्वयंसिद्धक आहे.

$n>0$ असेल, तर सिद्धतेचा शेवट निगमनाने
होतो. $n$ साठी दावा सिद्ध करायचा आहे; परंतु संदर्भ
वेगळा असला तरी उंची~$<n$ असलेल्या कोणत्याही
!!{proof} साठी दावा खरा आहे असे मानता येते. म्हणजे
कोणत्याही $k<n$ आणि कोणत्याही $\Pi$ व $\Lambda$
साठी $!A \land !B, \Pi \Sequent \Lambda$ ची
!!a{proof} !!{height}~$k$ असेल, तर $!A, !B, \Pi
\Sequent \Lambda$ ची !!a{proof} !!{height}~$\le k$
असलेली मिळते.

दोन प्रसंग आहेत: डावीकडील $!A \land
!B$ प्रमुख असतो किंवा नसतो. तो प्रमुख असेल तर
शेवटचे निगमन~\LeftR{\land} असते आणि तत्काळ
उप-!!{proof}~$\pi_1$ ची अंतिम क्रमवर्ती
$!A, !B, \Gamma \Sequent \Delta$ असते. या
प्रसंगात निष्कर्ष मिळतो, कारण $\pheight{\pi_1} <
\pheight{\pi}$.

$!A \land !B$ प्रमुख नसेल तर दुसरे
एखादे !!{formula} प्रमुख असते आणि $!A
\land !B$ हा शेवटच्या निगमनाच्या संदर्भात असतो.
मग शेवटच्या नियम~$R$ च्या तत्काळ उप-!!{proof}s
ना विगमनात्मक गृहीतक लावता येते. त्यातून~$\pi$
च्या तत्काळ उप-!!{proof}s च्या उंचींपेक्षा मोठी
!!{height} नसलेली !!a{proof}~$\pi_1'$ किंवा
दोन !!{proof}s~$\pi_1'$ आणि~$\pi_2'$ मिळतात.
$\pi_1'$ आणि~$\pi_2'$ यांच्या अंतिम क्रमवर्तींमध्ये
डावीकडील संदर्भात $!A \land !B$ ऐवजी $!A, !B$
येतात. त्यांना नियम~$R$ लावून सिद्धता~$\pi'$
मिळते. उदाहरणार्थ,~$\pi$ चा शेवट~\LeftR{\lif}
ने होत असेल तर:
\[
\AxiomC{}
\RightLabel{$\pi_1$}
\Deduce$!A \land !B, \Gamma' \fCenter \Delta, !C$
\AxiomC{}
\RightLabel{$\pi_2$}
\Deduce$!A \land !B, \Gamma', !D \fCenter \Delta$
\RightLabel{\LeftR{\lif}}
\BinaryInf$!A \land !B, \Gamma', !C \lif !D \fCenter \Delta$
\DisplayProof
\]
येथे डावीकडील $!C \lif !D$ प्रमुख असून डावीकडील
संदर्भ $!A \land !B, \Gamma'$ आहे (म्हणजे
$\Gamma = \Gamma', !C \lif !D$). विगमनात्मक
गृहीतकाने अनुक्रमे $!A, !B, \Gamma' \Sequent \Delta, !C$
आणि $!A, !B, \Gamma', !D \Sequent \Delta$
यांच्या !!{proof}s~$\pi_1'$
आणि~$\pi_2'$ मिळतात. या दोन क्रमवर्ती पुन्हा
\LeftR{\lif} चे पूर्वपक्ष बनून~$\pi'$ देतात:
\[
\AxiomC{}
\RightLabel{$\pi_1'$}
\Deduce$!A, !B, \Gamma' \fCenter \Delta, !C$
\AxiomC{}
\RightLabel{$\pi_2'$}
\Deduce$!A, !B, \Gamma', !D \fCenter \Delta$
\RightLabel{\LeftR{\lif}}
\BinaryInf$!A, !B, \Gamma', !C \lif !D \fCenter \Delta$
\DisplayProof
\]
आपल्याला पुढील मिळते:
\begin{align*}
\pheight{\pi} & = \max(\pheight{\pi_1}, \pheight{\pi_2}) + 1\\
\pheight{\pi'} & = \max(\pheight{\pi_1'}, \pheight{\pi_2'}) + 1
\intertext{व्याख्येनुसार; आणि}
\pheight{\pi_1'} & \le \pheight{\pi_1}\\
\pheight{\pi_2'} & \le \pheight{\pi_2}
\intertext{विगमन गृहीतकानुसार. म्हणून,}
\pheight{\pi'} & \le \pheight{\pi}.
\end{align*}
\end{proof}

\begin{prob}
\RightR{\lif} साठी \olref{lem:G3c-invert} ची सिद्धता द्या.
\end{prob}

\begin{lem}\ollabel{lem:invert-quant}
\RightR{\lforall} आणि \LeftR{\lexists} हे नियम
पुढील अर्थाने !!{height} जपून व्युत्क्रमणीय आहेत:
\begin{enumerate}
  \item $\Log{G3c} \Proves[n] \Gamma \Sequent \Delta,
  \lforall[x][!A(x)]$ असेल, तर $\Log{G3c} \Proves[n] \Gamma \Sequent
  \Delta, !A(c)$; आणि
  \item $\Log{G3c} \Proves[n] \lexists[x][!A(x)], \Gamma \Sequent
  \Delta$ असेल, तर $\Log{G3c} \Proves[n] !A(c), \Gamma \Sequent \Delta$,
\end{enumerate}
येथे प्रत्येक प्रसंगात~$c$ हा अनुक्रमे $\Gamma$, $\Delta$
आणि $\lforall[x][!A(x)]$ किंवा $\lexists[x][!A(x)]$
यांत न येणारा कोणताही स्थिरांक आहे.
\end{lem}

\begin{proof}
  आपण (1) सिद्ध करू आणि (2) प्रश्न म्हणून सोडू.
  दिलेल्या~$c$ पासून सर्व अंतर्गत आयगेन स्थिरांक
  वेगळे होतील अशी सुरक्षित नियमित नावे आधी घ्या.
  युक्तिवाद \olref{lem:G3c-invert} च्या सिद्धतेसारखाच
  आहे. विगमनात्मक टप्प्यात पुन्हा दोन प्रसंग असतात.
  $\lforall[x][!A(x)]$ प्रमुख असण्याचा प्रसंग सोपा आहे:
  शेवटच्या $\RightR{\lforall}$ निगमनाचा पूर्वपक्ष
  अपेक्षित $\Gamma
  \Sequent \Delta, !A(a)$ या रूपाचा असतो आणि
  त्यावर संपणाऱ्या तत्काळ उप-!!{proof}~$\pi_1$ ची
  $\pheight{\pi_1} < \pheight{\pi}$ असते.
  शिवाय, स्वचलाच्या अटीमुळे $a$ हा
  $\Gamma$, $\Delta$, किंवा $\lforall[x][!A(x)]$
  यांत येत नाही. $\pi_1$ नियमित आहे असे मानता
  येत असल्याने \olref[qua]{lem:subst-regular} नुसार
  $\Subst{\pi_1}{c}{a}$ ही $\Gamma \Sequent \Delta, !A(c)$
  ची त्याच !!{height} ची !!a{proof} आहे.
उंची~$0$ चा प्रसंगही जपला जातो: संख्यापकीय सूत्र
अणुरूप नसल्याने समान अणुरूप सूत्राचा किंवा डावीकडील
असत्याचा साक्षी उरलेल्या संदर्भात आहे; म्हणून त्याची
जागा~$!A(c)$ ने बदलल्यावरही क्रमवर्ती स्वयंसिद्धक आहे.
नव्या नियमाच्या आयगेन स्थिरांकांनाही~$c$ पासून वेगळे
नावे देऊन उरलेल्या विधानीय प्रसंगांप्रमाणे नियम पुन्हा
लावायचा.

दुसऱ्या प्रसंगात $\lforall[x][!A(x)]$ हे
  शेवटच्या निगमनात प्रमुख नसते; दुसरे सूत्र प्रमुख
  असते. पुन्हा एका उदाहरणाने प्रसंग पाहू. शेवटचे
  निगमन~\RightR{\land} आहे असे समजा.
  $\Delta$ मधील एखादे !!{formula}~$!B \land !C$
  या रूपाचे असून प्रमुख आहे. !!{proof}~$\pi$
  अशी संपते:
\[
\AxiomC{}
\RightLabel{$\pi_1$}
\Deduce$\Gamma \fCenter \Delta', \lforall[x][!A(x)], !B$
\AxiomC{}
\RightLabel{$\pi_2$}
\Deduce$\Gamma \fCenter \Delta', \lforall[x][!A(x)], !C$
\RightLabel{\RightR{\land}}
\BinaryInf$\Gamma \fCenter \Delta', \lforall[x][!A(x)], !B \land !C$
\DisplayProof
\]
येथे $\Delta = \Delta', !B \land !C$.
$\Delta$ मध्ये आणि पूर्ण निष्कर्षाच्या उरलेल्या
भागातही न येणारा !!a{constant}~$c$ घ्या;
तो~$\Delta', !B$ किंवा $\Delta', !C$ मध्येही
येत नाही हे उघड आहे. म्हणून विगमनात्मक गृहीतक
$\pi_1$ आणि~$\pi_2$ ना लागू होते; आपल्याकडे
!!{proof}s~$\pi_1'$ आणि $\pi_2'$ आहेत, ज्यांसाठी
$\pheight{\pi_1'} \le \pheight{\pi_1}$ आणि
$\pheight{\pi_2'} \le
\pheight{\pi_2}$. अपेक्षित $\Gamma \Sequent
\Delta, !A(c)$ ची !!{proof}~$\pi'$ अशी:
\[
\AxiomC{}
\RightLabel{$\pi_1'$}
\Deduce$\Gamma \fCenter \Delta', !A(c), !B$
\AxiomC{}
\RightLabel{$\pi_2'$}
\Deduce$\Gamma \fCenter \Delta', !A(c), !C$
\RightLabel{\RightR{\land}}
\BinaryInf$\Gamma \fCenter \Delta', !A(c), !B \land !C$
\DisplayProof
\]
आणि $\pheight{\pi'} = \max(\pheight{\pi_1'}, \pheight{\pi_2'})+1 \le \pheight{\pi}$.
\end{proof}

\Olref{lem:G3c-invert} चा कट-निरसन प्रमेयाच्या
सिद्धतेत महत्त्वाचा उपयोग होतो. पण त्याने पुढील
निष्कर्षही दाखवता येतो:

\begin{prop}\ollabel{prop:G3c-cont-adm}
\Log{G1c} चे आकुंचन नियम~\LeftR{\Contraction} आणि
\RightR{\Contraction} हे \Log{G3c} मध्ये
!!{height} जपून अनुज्ञेय आहेत.
\end{prop}

\begin{proof}
\RightR{\Contraction} आणि \LeftR{\Contraction}
दोन्हींसाठी एकाच वेळी युक्तिवाद देऊ.
\Log{G3c} मध्ये $\pi$ ही $\Gamma \Sequent
\Delta, !A, !A$ ची किंवा $!A, !A, \Gamma \Sequent
\Delta$ ची !!a{proof} आहे असे माना. अनुक्रमे
$\Gamma \Sequent \Delta, !A$ किंवा $!A,
\Gamma \Sequent \Delta$ ची अशी \Log{G3c}-!!{proof}
~$\pi'$ आहे की $\pheight{\pi'} \le
\pheight{\pi}$, हे दाखवायचे आहे. ते
$n = \pheight{\pi}$ वर विगमन करून करू.

$n=0$ असेल, तर~$\pi$ हे फक्त स्वयंसिद्धक
आहे. $\Gamma \Sequent \Delta, !A,
!A$ हे \Log{G3c} चे स्वयंसिद्धक असेल तर
$\Gamma \Sequent \Delta, !A$ हेही स्वयंसिद्धक
असते: कारण एखादा अण्वीय~$!C$ हा~$\Gamma$
आणि~$\Delta$ दोहोंचा !!a{element} असतो, किंवा
$!A$ अण्वीय असून~$\Gamma$ चा !!a{element}
असतो. अंतिम क्रमवर्ती $!A, !A, \Gamma \Sequent
\Delta$ असेल, तरी हाच युक्तिवाद लागू होतो.
डावीकडील असत्यामुळे स्वयंसिद्धक असेल, तर आकुंचनानंतर
किमान एक असत्य आस्थिती राहते; म्हणून तो प्रसंगही
उंची~$0$ चा असतो.

$n>0$ असेल, तर~$\pi$ चा शेवट काही
नियम~$R$ ने होतो. या शेवटच्या निगमनात $!A$
प्रमुख नसेल, तर \olref{lem:G3c-invert} च्या
सिद्धतेप्रमाणे~$\pi$ च्या तत्काळ उप-!!{proof}s
ना विगमनात्मक गृहीतक लावून आणि मिळालेल्या
!!{proof}(s) ना पुन्हा तोच नियम~$R$ लावून
!!a{proof}~$\pi'$ मिळवता येते. विगमनात्मक गृहीतक
असे आहे: $\Pi \Sequent \Lambda, !D,
!D$ किंवा $!D, !D, \Pi \Sequent \Lambda$ यांची
!!{height}~$k<n$ असलेली !!a{proof} असेल, तर
अनुक्रमे $\Pi \Sequent \Lambda, !D$ किंवा $!D, \Pi \Sequent
\Lambda$ यांची !!{height}~$\le k$ असलेली
!!a{proof} आहे. (संदर्भ किंवा आकुंचित होणारे
!!{formula} हे~$\Gamma$, $\Delta$, किंवा~$!A$
यांहून वेगळे असले तरी विगमनात्मक गृहीतक लागू होते!)

उदाहरणार्थ, $R$ हा~$\RightR{\land}$
असून~$\pi$ चा शेवट असा आहे असे समजा:
\[
\AxiomC{}
\RightLabel{$\pi_1$}
\Deduce$\Gamma \fCenter \Delta', !B, !A, !A$
\AxiomC{}
\RightLabel{$\pi_2$}
\Deduce$\Gamma \fCenter \Delta', !C, !A, !A$
\RightLabel{\RightR{\land}}
\BinaryInf$\Gamma \fCenter \Delta', !B \land !C, !A, !A$
\DisplayProof
\]
येथे $\Delta = \Delta', !B \land !C$. विगमनात्मक
गृहीतकाने अनुक्रमे $\Gamma \fCenter \Delta', !B, !A$
आणि $\Gamma \fCenter \Delta', !C, !A$ यांच्या
!!{proof}s~$\pi_1'$ आणि $\pi_2'$ मिळतात, आणि
$\pheight{\pi_1'} \le  \pheight{\pi_1}$ व
$\pheight{\pi_2'} \le
\pheight{\pi_2}$ असते. मग !!{proof}~$\pi'$ अशी:
\[
\AxiomC{}
\RightLabel{$\pi_1'$}
\Deduce$\Gamma \fCenter \Delta', !B, !A$
\AxiomC{}
\RightLabel{$\pi_2'$}
\Deduce$\Gamma \fCenter \Delta', !C, !A$
\RightLabel{\RightR{\land}}
\BinaryInf$\Gamma \fCenter \Delta', !B \land !C, !A$
\DisplayProof
\]

शेवटच्या निगमनात~$!A$ प्रमुख असेल, तर
~$\pi$ च्या तत्काळ उप-!!{proof}s ना व्युत्क्रमण
पूर्वप्रमेय (\olref{lem:G3c-invert}) लावतो, मग
विगमनात्मक गृहीतक लावतो, आणि शेवटी नियम~$R$
पुन्हा लावतो. उदाहरणार्थ, $\pi$ ही
$\Gamma \Sequent \Delta, !A, !A$ ची सिद्धता
असून $!A
\ident (!B \land !C)$ हे शेवटच्या निगमनात
प्रमुख आहे असे समजा. मग~$\pi$ चे रूप असे:
\[
\AxiomC{}
\RightLabel{$\pi_1$}
\Deduce$\Gamma \fCenter \Delta, !B \land !C, !B$
\AxiomC{}
\RightLabel{$\pi_2$}
\Deduce$\Gamma \fCenter \Delta, !B \land !C, !C$
\RightLabel{\RightR{\land}}
\BinaryInf$\Gamma \fCenter \Delta, !B \land !C, !B \land !C$
\DisplayProof
\]
\olref{lem:G3c-invert} हे~$\pi_1$ ला लावल्यास
$\Gamma \Sequent \Delta, !B, !B$ ची
!!a{proof}~$\pi_1'$ मिळते. (त्यातून
$\Gamma \Sequent \Delta, !C, !B$ ची !!a{proof}
सुद्धा मिळते, पण आपल्याला ती लागत नाही.)
तसेच~$\pi_2$ ला \olref{lem:G3c-invert} लावल्यास
$\Gamma \Sequent \Delta, !C, !C$ ची
!!a{proof}~$\pi_2'$ मिळते. \RightR{\land} चे
व्युत्क्रमण !!{height} जपते, म्हणून
$\pheight{\pi_1'} \le  \pheight{\pi_1} < \pheight{\pi}$
आणि $\pheight{\pi_2'} \le  \pheight{\pi_2} < \pheight{\pi}$.
त्यामुळे~$\pi_1'$ आणि $\pi_2'$ ना विगमनात्मक गृहीतक
लागू होते: अनुक्रमे $\Gamma \Sequent \Delta, !B$
आणि $\Gamma \Sequent \Delta, !C$ यांच्या
!!{proof}s~$\pi_1''$ आणि $\pi_2''$ मिळतात, आणि
$\pheight{\pi_1''}
\le \pheight{\pi_1'} \le \pheight{\pi_1}$ तसेच
$\pheight{\pi_2''} \le
\pheight{\pi_2'} \le \pheight{\pi_2}$ असते.
मग !!{proof}~$\pi'$ अशी:
\[
\AxiomC{}
\RightLabel{$\pi_1''$}
\Deduce$\Gamma \fCenter \Delta, !B$
\AxiomC{}
\RightLabel{$\pi_2''$}
\Deduce$\Gamma \fCenter \Delta, !C$
\RightLabel{\RightR{\land}}
\BinaryInf$\Gamma \fCenter \Delta, !B \land !C$
\DisplayProof
\]
तिच्यासाठी $\pheight{\pi'} \le \pheight{\pi}$.


!!{formula} प्रमुख असताना संख्यकांच्या
नियमांबाबत विशेष काळजी घ्यावी लागते.

$!A \ident \lforall[x][!B(x)]$ हे उजवीकडे
प्रमुख आहे असे समजा. मग~$\pi$ चा शेवट असा:
\[
\AxiomC{}
\RightLabel{$\pi_1$}
\Deduce$\Gamma \fCenter \Delta, \lforall[x][!B(x)], !B(a)$
\RightLabel{\RightR{\lforall}}
\UnaryInf$\Gamma \fCenter \Delta, \lforall[x][!B(x)], \lforall[x][!B(x)]$
\DisplayProof
\]
\olref{lem:invert-quant} नुसार, $\Gamma \fCenter \Delta,
\lforall[x][!B(x)], !B(a)$ मध्ये न येणाऱ्या
कोणत्याही~$c$ साठी $\Gamma \fCenter \Delta, !B(c),
!B(a)$ ची !!{height}~$\le \pheight{\pi_1}$
असलेली !!a{proof}~$\pi_1'$ मिळते. $\pi_1'$
नियमित करून तिचे अंतर्गत आयगेन स्थिरांक~$a$ पासून
वेगळे आहेत असे मानता येते; म्हणून
\olref[qua]{lem:subst-regular} नुसार
!!{proof}~$\Subst{\pi_1'}{a}{c}$ ही $\Gamma \fCenter \Delta, !B(a),
!B(a)$ ची त्याचप्रमाणे !!{height}~$\le \pheight{\pi_1}$
असलेली !!a{proof} आहे. आता विगमनात्मक गृहीतक
लागू होऊन $\Gamma
\fCenter \Delta, !B(a)$ ची !!a{proof}~$\pi_1''$
मिळते. \RightR{\lforall} नियम लावून~$\pi'$
मिळवू:
\[
\AxiomC{}
\RightLabel{$\pi_1''$}
\Deduce$\Gamma \fCenter \Delta, !B(a)$
\RightLabel{\RightR{\lforall}}
\UnaryInf$\Gamma \fCenter \Delta, \lforall[x][!B(x)]$
\DisplayProof
\]
आणि $\pheight{\pi'} \le \pheight{\pi}$.


आता $!A \ident \lexists[x][!B(x)]$ हे
उजवीकडे प्रमुख आहे असे समजा. मग~$\pi$ चे रूप असे:
\[
\AxiomC{}
\RightLabel{$\pi_1$}
\Deduce$\Gamma \fCenter \Delta, \lexists[x][!B(x)], \lexists[x][!B(x)], !B(t)$
\RightLabel{\RightR{\lexists}}
\UnaryInf$\Gamma \fCenter \Delta, \lexists[x][!B(x)], \lexists[x][!B(x)]$
\DisplayProof
\]
विगमनात्मक गृहीतक~$\pi_1$ ला लागू होते; त्यामुळे
$\Gamma \Sequent \Delta, \lexists[x][!B(x)], !B(t)$
ची !!a{proof}~$\pi_1'$ मिळते.
(येथे व्युत्क्रमण पूर्वप्रमेयाची गरज नाही!) मग
!!{proof}~$\pi'$ अशी:
\[
\AxiomC{}
\RightLabel{$\pi_1'$}
\Deduce$\Gamma \fCenter \Delta,  \lexists[x][!B(x)], !B(t)$
\RightLabel{\RightR{\lexists}}
\UnaryInf$\Gamma \fCenter \Delta, \lexists[x][!B(x)]$
\DisplayProof
\]
तिच्यासाठी $\pheight{\pi'} \le \pheight{\pi}$.
\end{proof}

\begin{prob}
\LeftR{\Contraction} साठी \olref{prop:G3c-cont-adm}
च्या सिद्धतेची रूपरेषा द्या. $!A \ident (!B \land !C)$
आणि $!A \ident (!B
\lif !C)$ हे प्रसंग वर्णा.
\end{prob}

\begin{prob}
उरलेल्या संख्यकांच्या प्रसंगांसाठी, म्हणजे $!A \ident
\lforall[x][!B(x)]$ आणि $!A \ident \lexists[x][!B(x)]$
असताना, \LeftR{\Contraction} साठी
\olref{prop:G3c-cont-adm} च्या सिद्धतेची रूपरेषा द्या.
\end{prob}

\end{document}
